\chapter{Epilogue}
\label{chap:epilogue}

Human DNA consists of about three billion base pairs, and our brains contain roughly one hundred billion neurons.
The number of interconnections between these neurons is estimated at around 100 trillion---far more complex and deeply integrated than the largest artificial neural networks in existence today.

Given the rapid pace of technological advancement, it is conceivable that artificial neural networks may one day surpass the human brain in sheer architectural complexity.
Perhaps this will happen within the next decade or two.
Yet, even if they do, raw complexity alone may not capture the full richness of the human experience.
Time will tell.

Static highly compressed set of information doesn't allow structures in it - us - to create new information.
Correspondingly, we will never be able to create a new conscious human as a computer simulation,
but explore those that already exist.
And even if we could, simulating a single human brain at the level of every individual neuron would require an astronomical amount of computing power---a network of hundreds of billions of CPUs, occupying cubic kilometers of space and drawing the energy of a nuclear power plant.
Furthermore, to simulate a whole person convincingly, one would need to simulate a universe around them, complete with trillions of cells and the untold billions of stars of a vast, surrounding galaxy.

In short, God---if there is one---does not need to worry about simulated souls any time soon.
It is surely we humans who have some catching up to do.

In \textit{A Brief History of Time}, Stephen Hawking famously asked what ``breathes fire into the equations and makes a universe for them to describe.''
The traditional scientific approach---constructing mathematical models---cannot answer why there should be a universe at all.
Why bother existing?

The conclusion proposed in this book is that nothing breathes fire into anything.
There is no ``universe'' beyond the abstract equations we call the laws of physics.
The nature of everything is fundamentally abstract.

At one point, I considered naming this book something more serious, like \textit{The Real Theory of Everything}.
However, exercising my illusion of free will, I chose otherwise.
Besides, that name was already taken.

It is so easy to imagine total nothingness---the absolute absence of matter and energy---as the default state of reality.
This intuition arises from our evolutionary hardwiring: survival requires ``having something''---air to breathe, solid ground beneath our feet, food to eat.
If we lack these, we die.
Yet, when considering the mathematical set of all possible sets, it is statistically probable that there is something rather than nothing.

And despite being nothing but abstract information, it feels astonishingly real.
